\documentclass[11pt]{article}
\usepackage[margin=0.9in]{geometry}
\usepackage{amsmath,amssymb,mathtools,bm}
\usepackage{microtype}
\usepackage{hyperref}
\usepackage{enumitem}
\newcommand{\bg}{\bar g}
\newcommand{\bn}{\bar\nabla}
\newcommand{\be}{\bar e}
\newcommand{\bt}{\bar\theta}
\newcommand{\bw}{\bar\omega}
\newcommand{\bR}{\bar R}
\newcommand{\cC}{\mathcal C}
\newcommand{\cE}{\mathcal E}
\newcommand{\Lie}{\mathcal L}
\newcommand{\dd}{\mathrm d}
\newcommand{\order}{\mathcal O}
\title{A Symmetric-Hyperbolic Reference-Generalized-Harmonic Formulation\\for Dynamical Wormhole-to-Trumpet Puncture Evolution}
\author{Theory/static-code audit of AthenaK Ref-GH}
\date{31 August 2026}
\begin{document}
\maketitle

\begin{abstract}
We derive and audit a first-order generalized-harmonic (GH) formulation intended for no-excision puncture evolution beginning with ordinary Bowen--York wormhole data and dynamically approaching a moving-puncture trumpet. The existing STANDARD $\gamma_1=-1$ GH principal part is symmetric hyperbolic, but symmetric hyperbolicity does not prevent unstable lower-order modes generated by a strongly varying reference frame. We show that a time-dependent singular exponent $r^{-q(t)}$ is unsuitable because $\partial_t\ln r^{-q(t)}=-\dot q\ln r$ and related mixed derivatives are unbounded. We recommend instead evolving the deviation $h_{AB}=\Psi_{AB}-\eta_{AB}$ and its background-covariant first derivatives, retaining an algebraic wave-map plus relative-damped gauge, and carrying the Bowen--York wormhole-to-trumpet coordinate transition in a genuine dynamical single-hole reference trajectory. The exact puncture remains on the wormhole branch at every finite coordinate time while every fixed $r>0$ approaches the trumpet. The cleanest single-hole reference is a Schwarzschild gauge-transition solution, for which the zero-residual state is exact. We give the characteristic fields, a positive symmetrizer, covariant reduction/curl constraints, puncture regularity conditions, and a GH-consistent initial-data construction.
\end{abstract}

\section{Scope and final decision}
The target is
\[
\boxed{\text{first order}+\text{symmetric hyperbolic}+\text{puncture compatible}+\text{dynamically gauge controlled}}
\]
without excision and without relying on cancellation between separately divergent quantities.

The recommended system is STANDARD $\gamma_1=-1$ first-order GH in reference-deviation variables and reference-covariant first derivatives, with a wave-map plus relative-damped algebraic gauge and a genuine dynamical wormhole-to-trumpet reference trajectory. The formulation is explicitly intended to start from ordinary isotropic Bowen--York wormhole data; only the scalar dynamic-exponent representation is rejected.

\section{Reduced GH equations}
Let $g_{ab}$ have signature $(-,+,+,+)$ and define
\[
C_a=H_a+\Gamma_a,\qquad \Gamma_a=g^{bc}\Gamma_{abc}.
\]
Use the damped reduced Einstein equation
\begin{equation}
R_{ab}-\nabla_{(a}C_{b)}+\frac{\gamma_0}{2}
\left(2n_{(a}C_{b)}-g_{ab}n^cC_c\right)=0.
\label{eq:reduced}
\end{equation}
Let $\bg_{ab}$ be a prescribed reference metric with orthonormal frame $\be_A{}^a$, coframe $\bt^A{}_a$, spin connection $\bw^A{}_{BC}$, structure coefficients $c^A{}_{BC}$, curvature $\bR^A{}_{BCD}$, and derivative $\bn$. Define
\[
\Psi_{AB}=\be_A{}^a\be_B{}^b g_{ab}.
\]
At exact match, $g=\bg$ and $\Psi_{AB}=\eta_{AB}$.

\section{Audit of the STANDARD principal part}
For one symmetric component, the old first-order variables are
\[
\Psi,\qquad \Pi=-n^a\partial_a\Psi,\qquad \Phi_I=E_I\Psi.
\]
The anholonomic curl is
\[
C_{IJ}=E_I\Phi_J-E_J\Phi_I-c^K{}_{IJ}\Phi_K.
\]
The STANDARD ordering is
\begin{equation}
(\partial_t\Phi_I)_{\rm std}=(\partial_t\Phi_I)_{\rm comp}-\beta^J C_{IJ},
\label{eq:standard}
\end{equation}
which is the sign implemented by the audited source.

For $\gamma_1=-1$, the standard reduction-damping additions are
\[
\delta(\partial_t\Pi)=-\gamma_2\beta^iC_i,
\qquad
\delta(\partial_t\Phi_I)=\alpha\gamma_2\be_I{}^i C_i.
\]
Let $s_I$ satisfy $G^{IJ}s_Is_J=1$. The characteristic fields are
\begin{align}
u^0&=\Psi,\\
u^\pm&=\Pi\pm s^I\Phi_I-\gamma_2\Psi,\\
u_I^\perp&=\Phi_I-s_I s^J\Phi_J,
\end{align}
with speeds
\[
0,\qquad -\beta^s+\alpha,\qquad -\beta^s-\alpha,\qquad -\beta^s.
\]
A one-component symmetrizer block is
\begin{equation}
S=\begin{pmatrix}
\Lambda^2&-\gamma_2&0\\
-\gamma_2&1&0\\
0&0&1
\end{pmatrix},
\label{eq:Sblock}
\end{equation}
with the $Q$ block generalized by $G^{IJ}$. Positivity follows from
\[
P^2-2\gamma_2hP+\Lambda^2h^2=(P-\gamma_2h)^2+(\Lambda^2-\gamma_2^2)h^2,
\]
so $\Lambda^2>\gamma_2^2$ is sufficient. Exact symbolic multiplication verifies $SA^i=(SA^i)^T$.

\section{Why a fixed reference can still grow}
For
\[
\partial_tu+A^i(x)\partial_i u=B(x)u,\qquad SA^i=(SA^i)^T,
\]
the energy estimate contains
\begin{equation}
\frac{\dd E}{\dd t}\le
\int u^T\left[SB+B^TS+\partial_i(SA^i)\right]u\,\dd^3x+\text{boundary flux}.
\label{eq:energy}
\end{equation}
Thus symmetric hyperbolicity does not require the lower-order operator to be dissipative. Spatial-frame gradients, structure coefficients, spin terms, reference curvature, and the gauge Jacobian can generate a positive local energy-growth matrix even when the principal symbol is perfectly symmetric hyperbolic. This explains how the fixed-reference reduction/curl mode can be structurally real without being a principal-symbol defect.

\section{Recommended evolved variables}
Define
\begin{equation}
\boxed{h_{AB}=\Psi_{AB}-\eta_{AB}.}
\end{equation}
Evolve background-covariant first derivatives
\begin{equation}
\boxed{P_{AB}=-n^a\bn_a h_{AB},\qquad Q_{IAB}=\be_I{}^a\bn_a h_{AB}.}
\label{eq:variables}
\end{equation}
Since $\bn\eta=0$, these are also covariant derivatives of $\Psi$. In frame components,
\begin{equation}
Q_{IAB}=E_Ih_{AB}-\bw^C{}_{AI}h_{CB}-\bw^C{}_{BI}h_{AC}.
\label{eq:Qdef}
\end{equation}
The exact matched state is
\[
\boxed{h=P=Q=0.}
\]
The connection difference is directly
\begin{equation}
\Delta_{ABC}=\frac12\left(Q_{BAC}+Q_{CAB}-Q_{ABC}\right).
\label{eq:Delta}
\end{equation}
This eliminates the raw-derivative minus spin-connection cancellation used to reconstruct $Q$ in the current production path.

For every $r>0$, the change from $(\Psi,\Pi,\Phi)$ to $(h,P,Q)$ is affine and triangular with identity derivative-variable diagonal blocks, so only lower-order terms change. The new system is nevertheless regarded as native at the puncture; no assumption of a uniformly bounded old-to-new transformation at $r=0$ is needed.

\section{First-order evolution and symmetrizer}
The principal part is
\begin{align}
\partial_t h_{AB}&\simeq-\alpha P_{AB}+\beta^I Q_{IAB},
\label{eq:hprincipal}\\
\partial_tP_{AB}&\simeq\beta^IE_IP_{AB}-\alpha G^{IJ}E_IQ_{JAB}
-\gamma_2\beta^I\cC_{IAB},
\label{eq:Pprincipal}\\
\partial_tQ_{IAB}&\simeq\beta^JE_JQ_{IAB}-\alpha E_IP_{AB}
+\alpha\gamma_2\cC_{IAB}.
\label{eq:Qprincipal}
\end{align}
All spin, curvature, reference-motion, lapse/shift-gradient, and algebraic-gauge terms are lower order and must be written directly in background-covariant residual form.

For positive-definite $m^{AB}$, define
\begin{equation}
\begin{split}
\cE={}&m^{AC}m^{BD}\bigl[\Lambda^2h_{AB}h_{CD}+P_{AB}P_{CD}
-2\gamma_2h_{AB}P_{CD}\\
&\hspace{38mm}+G^{IJ}Q_{IAB}Q_{JCD}\bigr].
\end{split}
\label{eq:energynew}
\end{equation}
If $G^{IJ}>0$ and $\Lambda^2>\gamma_2^2$, Eq.~\eqref{eq:energynew} is positive definite and symmetrizes Eqs.~\eqref{eq:hprincipal}--\eqref{eq:Qprincipal}. Thus the proposed main system is genuinely symmetric hyperbolic.

\section{Covariant subsidiary constraints}
Define the reduction constraint
\begin{equation}
\boxed{\cC_{IAB}=\bn_Ih_{AB}-Q_{IAB}.}
\label{eq:redC}
\end{equation}
For an anholonomic frame define
\begin{equation}
\boxed{\begin{split}
\cC_{IJAB}={}&\bn_IQ_{JAB}-\bn_JQ_{IAB}-c^K{}_{IJ}Q_{KAB}\\
&+\bR^C{}_{A IJ}h_{CB}+\bR^C{}_{B IJ}h_{AC}.
\end{split}}
\label{eq:curlC}
\end{equation}
The sign is fixed so that Eq.~\eqref{eq:curlC} vanishes identically when Eq.~\eqref{eq:redC} holds, using the anholonomic commutator of $\bn$.

The principal reduction-constraint propagation is
\begin{equation}
\partial_t\cC_{IAB}\simeq-\alpha\gamma_2\cC_{IAB}+\beta^J\cC_{JIAB},
\end{equation}
with remaining terms algebraic in the reduction/curl constraints and bounded reference curvature/connection coefficients. The curl is transported by the shift and coupled algebraically to reduction errors and curvature.

\section{Regular GH constraint and algebraic gauge}
Take the wave-map base source
\[
B_a=-g_{ab}g^{cd}\bar\Gamma^b{}_{cd}
\]
and write $H_a=B_a+J_a$. Then
\[
C^a=g^{bc}(\Gamma^a{}_{bc}-\bar\Gamma^a{}_{bc})+J^a,
\]
so in frame lower-index form
\begin{equation}
\boxed{C_A=\Delta_A+J_A.}
\label{eq:regularC}
\end{equation}
This is the quantity that should be evaluated in production; one should not separately form large $H_A$ and $\Gamma_A$ and subtract them.

Retain the relative-damped gauge
\begin{equation}
\boxed{H_a=B_a+\bt^A{}_aW D_A,\qquad
D_A=\mu_L L_RN_A^R-\frac{\mu_S}{a_R}V_A^R.}
\label{eq:gauge}
\end{equation}
The relative quantities are algebraic functions of $\Psi=\eta+h$. Use an exact puncture plateau $W=0$ in an open core, so the core is pure wave-map to the dynamical reference.

At exact match, $a_R=1$, $L_R=0$, and $V_A^R=0$, hence $D_A=J_A=0$ exactly.

\subsection{Linearized gauge modes}
At first order about $\Psi=\eta$,
\begin{equation}
D_0=-\frac{\mu_L}{2}(h_{00}+h_{ii}),\qquad D_i=-\mu_Sh_{0i}.
\end{equation}
For a pure gauge perturbation $h_{ab}=2\partial_{(a}\xi_{b)}$, transverse modes obey
\[
s^2+\mu_Ss+k^2=0.
\]
The scalar plane-wave matrix is
\[
M=\begin{pmatrix}
A+\mu_Ls&i\mu_Lk\\
i\mu_Sk&A+\mu_Ss
\end{pmatrix},\qquad A=s^2+k^2,
\]
and exact symbolic factorization gives
\begin{equation}
\boxed{\det M=(s^2+k^2)\left[s^2+(\mu_L+\mu_S)s+k^2+\mu_L\mu_S\right].}
\label{eq:gaugefactor}
\end{equation}
For $\mu_L,\mu_S>0$ there is no frozen-coefficient exponentially growing gauge root. One scalar harmonic mode remains undamped. Since the gauge is algebraic in $h$ and the prescribed reference, it does not alter the first-order principal symbol.

\section{No-go theorem for a time-dependent puncture exponent}
Suppose a reference scale contains $L_q(r,t)=r^{-q(t)}$. Then exactly
\begin{equation}
\boxed{\partial_t\ln L_q=-\dot q\ln r.}
\label{eq:qlog}
\end{equation}
Higher derivatives of $L_q$ contain $\ddot q\ln r$, $\dot q^2(\ln r)^2$, and mixed derivatives contain $\dot q/r$. Therefore a finite-time global change of the leading puncture exponent cannot be represented by a smooth scalar $q(t)$ while keeping production coefficients bounded. Dynamic `q` tracking should be abandoned in the puncture production path.

\section{Bowen--York to trumpet asymptotics}
For ordinary Bowen--York wormhole data,
\[
\psi_{\rm W}=1+\frac{M}{2r}+u,\qquad u=\order(1),
\]
so with $L=\psi^2$,
\[
L_{\rm W}\sim r^{-2}.
\]
A precollapsed lapse obeys $\alpha_{\rm W}\sim r^2$ and the standard initial shift is zero. Finite Bowen--York momentum and spin do not alter the leading conformal-factor power.

A stationary moving-puncture trumpet has finite limiting areal radius $R_0$, hence
\[
L_{\rm T}=\frac{R}{r}\sim r^{-1}.
\]
Its lapse vanishes as a positive power and its Cartesian shift is linear in $r$ at leading order.

The transition is therefore nonuniform at the mathematical puncture. A bounded no-excision formulation should keep the exact puncture on the wormhole branch at every finite coordinate time while a transition region moves inward; for every fixed $r>0$, the late-time geometry approaches the trumpet. This is the appropriate continuum representation of moving-puncture wormhole-to-trumpet relaxation.

\section{Preferred dynamical reference}
Construct in spherical symmetry a reference spacetime $\bg_{ab}(t,r)$ that is Schwarzschild written in coordinates satisfying:
\begin{enumerate}
\item at $t=0$, the selected isotropic wormhole slice, initial lapse convention, and zero shift;
\item during the transition, a regular moving-puncture-type coordinate evolution;
\item as $t\to\infty$, the stationary $1+\log$ trumpet profiles used by the existing table.
\end{enumerate}
This may be generated offline by a high-accuracy one-dimensional coordinate/gauge evolution on exact Schwarzschild. Production should store either the coordinate map or frame-native reference data with consistent first and second jets.

The preferred single-hole reference satisfies
\begin{equation}
\boxed{\bar R_{ab}=0}
\end{equation}
and has bounded frame-native coefficients
\begin{equation}
|\bw^A{}_{BC}|\le C_1/M,\qquad |\bR^A{}_{BCD}|\le C_2/M^2
\label{eq:refbounds}
\end{equation}
with finite puncture limits for every coefficient actually evaluated by the residual equations.

For this reference,
\begin{equation}
\boxed{h=P=Q=0,\qquad J=0,\qquad C_A=0}
\end{equation}
is the exact single-hole dynamical transition. Thus a large physical gauge change becomes a stationary zero residual, providing the strongest possible regression oracle.

An arbitrary smooth interpolation between wormhole and trumpet profiles is second best because it is generally nonvacuum and drives a finite residual source even at $g=\bg$.

\subsection{Reference coefficients must be regular before evaluation}
Do not compute a finite spin connection by subtracting separately divergent coordinate-frame terms such as $\partial e+\bar\Gamma e$ near the puncture. The reference provider should return regular frame-native connection coefficients directly, for example from orthonormal extrinsic curvature, spatial acceleration, triad rotation, and analytic Schwarzschild curvature. The same rule applies to reference curvature. A quantity that becomes finite only after subtracting divergent floating-point numbers fails the formulation gate.

\section{$C^\infty$ finite-radius plateaus}
For smooth finite-radius gauge/reference stitching define
\begin{equation}
\rho(x)=\begin{cases}0,&x\le0,\\ e^{-1/x},&x>0,\end{cases}
\qquad
S(x)=\frac{\rho(x)}{\rho(x)+\rho(1-x)}.
\label{eq:bump}
\end{equation}
Then $S=0$ for $x\le0$, $S=1$ for $x\ge1$, and $S\in C^\infty$ with all derivatives zero at the plateaus. For $0<x<1$, let
\[
L(x)=-\frac1x+\frac1{1-x},\qquad S=(1+e^{-L})^{-1}.
\]
Then
\begin{align}
S'&=S(1-S)L',&L'&=\frac1{x^2}+\frac1{(1-x)^2},\\
S''&=S(1-S)\left[(1-2S)(L')^2+L''\right],&L''&=-\frac2{x^3}+\frac2{(1-x)^3}.
\end{align}
For physical width $\Delta$, derivatives scale as $\Delta^{-1}$ and $\Delta^{-2}$. Smoothness alone is insufficient: $\Delta$ must remain finite in physical units and must not shrink with resolution or puncture-core radius.

\section{GH-consistent Bowen--York initial data}
Given physical $(\gamma_{ij},K_{ij})$, choose initial $\alpha$ and $\beta^i$ and construct
\begin{align}
g_{00}&=-\alpha^2+\gamma_{ij}\beta^i\beta^j,\\
g_{0i}&=\gamma_{ij}\beta^j,\\
g_{ij}&=\gamma_{ij}.
\end{align}
Compute the spatial first jet, then $h$ and $Q$ from Eqs.~\eqref{eq:variables}--\eqref{eq:Qdef}.

The ADM identity
\begin{equation}
\partial_t\gamma_{ij}=-2\alpha K_{ij}+\Lie_\beta\gamma_{ij}
\label{eq:K}
\end{equation}
fixes six independent time-derivative combinations. Use
\[
x^A=(\partial_t\alpha,\partial_t\beta^1,\partial_t\beta^2,\partial_t\beta^3)
\]
as the remaining four unknowns. For fixed Cauchy data and spatial derivatives,
\[
\Gamma_a=\Gamma_a^{(0)}+M_{aA}x^A.
\]
Since the recommended $H_a$ is algebraic in $g$ and the prescribed reference, solve
\begin{equation}
\boxed{M_{aA}x^A=-H_a-\Gamma_a^{(0)}}
\label{eq:4x4}
\end{equation}
to enforce $C_a=0$ without changing $(\gamma_{ij},K_{ij})$. The complete coordinate first jet is then known, and all ten independent components
\[
P_{AB}=-n^a\bn_a h_{AB}
\]
follow. Thus six combinations are fixed by $K_{ij}$ and four by the GH gauge constraint.

The reference should capture only the universal puncture singular factors. The regular Bowen--York correction $u$, momentum, and spin remain finite residual data in $h,P,Q$.

\section{Puncture regularity gate}
A production reference is admissible only if the coefficients actually appearing in the residual system have finite puncture limits. Require at minimum
\[
G^{IJ}=\order(1),\qquad (G^{-1})_{IJ}=\order(1),
\]
\[
\bw^A{}_{BC}=\order(M^{-1}),\qquad \bR^A{}_{BCD}=\order(M^{-2}),
\]
\[
J_A=\order(1),\qquad \frac{\partial J_A}{\partial h_{BC}}=\order(M^{-1}),
\]
with bounded derivatives required by the lower-order source. These bounds must hold before floating-point evaluation. The compactified puncture is a limiting point of the $r>0$ equations; no finite neighborhood is excised.

\section{Lower-order static qualification}
In addition to the principal-symbol proof, form the frozen-coefficient energy-growth matrix
\begin{equation}
K=S^{-1/2}\left[SB+B^TS+\partial_i(SA^i)\right]S^{-1/2}.
\label{eq:lowergate}
\end{equation}
Its largest eigenvalue bounds local energy growth. Apply the same construction to the subsidiary $(C_A,\cC_I,\cC_{IJ})$ system. This turns ``smooth enough reference'' into an auditable static condition and directly targets the observed fixed-reference mode.

\section{Spinning punctures and binaries}
Bowen--York boost and spin leave the leading wormhole conformal-factor power unchanged, so they enter as finite nonspherical residuals in $h,P,Q$. A spinning trumpet reference may later improve conditioning but is not required for the basic singular factorization.

For a binary, use two persistent local puncture cores, each with its own local wormhole-to-trumpet history and center. Stitch them to an outer binary frame only at finite radius with $C^\infty$ windows of fixed physical width. At common-horizon formation, keep both singular cores and transition only the smooth outer map/reference toward a remnant-centered frame. No abrupt two-singularity to one-singularity operation is required.

\section{Concrete AthenaK changes}
Retain STANDARD ordering, finite $\gamma_0$ and $\gamma_2$, the geometric connection-difference source, the relative-damped gauge idea, and the reference-provider abstraction. Deprecate dynamical `reference_q_controlled` exponent tracking, the live `Hhat/theta/Upsilon` moving-puncture driver in puncture production, separately divergent subtraction paths, and $C^2$ high-order production plateaus.

Concretely:
\begin{enumerate}
\item reinterpret the 10+10+30 Einstein fields as $h,P,Q$;
\item derive the STANDARD $Q$ equation directly rather than evolving raw $\Phi$;
\item feed fundamental $Q$ into the covariant source and eliminate the raw-to-covariant `frame_correction` path after oracle equivalence is established;
\item evaluate $C_A=J_A+\Delta_A$ directly;
\item add a `wormhole_trumpet_transition` reference provider with a consistent time-dependent two-jet and regular frame-native spin/curvature;
\item retain `relative_damped_gauge.hpp`, but use covariant first jets and a $C^\infty$ plateau;
\item add a reusable GH initial-data $4\times4$ solve implementing Eq.~\eqref{eq:4x4};
\item add static symmetrizer, coefficient-bound, and lower-order Jacobian gates.
\end{enumerate}

\section{Claim status}
\paragraph{Proved.}
The STANDARD $\gamma_1=-1$ principal symbol, characteristic speeds, positive symmetrizer, $\gamma_2$ sign, the no-go identity Eq.~\eqref{eq:qlog}, the relative-gauge factorization Eq.~\eqref{eq:gaugefactor}, and preservation of the principal symbol by the covariant derivative variables and algebraic gauge.

\paragraph{Supported by existing numerical evidence.}
The old live moving-puncture GH driver is unsafe in the present puncture realization; reference motion amplifies the instability; and a fixed strongly varying reference can support a common GH/reduction/curl lower-order growth mode.

\paragraph{Plausible but unproved.}
The covariant-derivative residual variables remove the observed fixed-reference growth mechanism; the dynamic Schwarzschild reference plus relative damping robustly controls generic Bowen--York single holes; and the two-core architecture remains robust through merger.

\paragraph{Requires future numerical testing.}
Production robustness, quantitative window widths, reference-table accuracy, and boosted/spinning/binary qualification.

\section{Exact symbolic checks}
For one spatial direction, the derivative matrix for $U=(h,P,Q)^T$ can be written
\[
B=\begin{pmatrix}
0&0&0\\
-\gamma_2\beta&\beta&-\alpha\\
\alpha\gamma_2&-\alpha&\beta
\end{pmatrix}.
\]
Together with Eq.~\eqref{eq:Sblock}, exact symbolic multiplication gives
\[
SB-(SB)^T=0.
\]
For the scalar relative-gauge matrix used above, exact factorization gives Eq.~\eqref{eq:gaugefactor}. These are exact algebraic identities rather than floating-point fits.

\section{Boxed final specification}
\[
\boxed{\begin{minipage}{0.94\linewidth}
\textbf{Evolved variables:} $U=(h_{AB},P_{AB},Q_{IAB})$ with Eq.~\eqref{eq:variables}.\\[0.4em]
\textbf{Reference:} a prescribed dynamical wormhole-to-trumpet reference $\bg_{ab}(t,x)$ with bounded frame-native connection/curvature coefficients; for the single-hole qualification case, a Schwarzschild vacuum gauge-transition solution.\\[0.4em]
\textbf{Gauge:} $H_a=-g_{ab}g^{cd}\bar\Gamma^b{}_{cd}+\bt^A{}_aWD_A$, with Eq.~\eqref{eq:gauge}, $\mu_L,\mu_S>0$, and an exact $C^\infty$ puncture plateau $W=0$.\\[0.4em]
\textbf{Evolution:} STANDARD $\gamma_1=-1$ GH principal part, Eqs.~\eqref{eq:hprincipal}--\eqref{eq:Qprincipal}, all lower-order terms written directly in background-covariant residual form.\\[0.4em]
\textbf{Constraints:} $C_A=J_A+\Delta_A$, Eq.~\eqref{eq:redC}, and the curvature-corrected curl Eq.~\eqref{eq:curlC}.\\[0.4em]
\textbf{Damping:} finite $\gamma_0>0$, finite $\gamma_2\ge0$; no lapse-divergent rescaling.\\[0.4em]
\textbf{Puncture regularization:} universal Bowen--York singular geometry carried by the reference; finite residual $h,P,Q$; no finite-time $q(t)$ exponent change; no floors/clipping/divergent subtraction.\\[0.4em]
\textbf{Reference control:} the exact puncture remains wormhole-like for every finite coordinate time while a regular transition region moves inward; every fixed $r>0$ approaches the stationary trumpet. For binaries, retain two local cores through merger and modify only smooth finite-radius outer maps.
\end{minipage}}
\]

\section*{References}
\begin{enumerate}[label={[rabic*]}]
\item L. Lindblom, M. A. Scheel, L. E. Kidder, R. Owen, O. Rinne, ``A New Generalized Harmonic Evolution System,'' arXiv:gr-qc/0512093.
\item L. Lindblom, B. Szilagyi, ``An Improved Gauge Driver for the Generalized Harmonic Einstein System,'' arXiv:0904.4873.
\item M. A. Scheel et al., ``Solving Einstein's Equations With Dual Coordinate Frames,'' arXiv:gr-qc/0607056, Phys. Rev. D 74, 104006.
\item M. Hannam et al., ``Geometry and Regularity of Moving Punctures,'' arXiv:gr-qc/0606099.
\item M. Hannam et al., ``Where do moving punctures go?,'' arXiv:gr-qc/0612097.
\item M. Hannam, S. Husa, N. O'Murchadha, ``Bowen--York trumpet data and black-hole simulations,'' arXiv:0908.1063.
\item T. Dietrich, B. Bruegmann, ``Solving the Hamiltonian constraint for 1+log trumpets,'' arXiv:1309.3087.
\end{enumerate}
\end{document}
